\documentclass{article}
\usepackage[utf8]{inputenc}
\usepackage[T1]{fontenc}
\usepackage{geometry}
\geometry{
    a4paper,
    left=2.5cm,
    right=2.5cm,
    top=2.5cm,
    bottom=2.5cm
}
\usepackage{graphicx}
\usepackage{float}
\usepackage{hyperref}
\usepackage[english]{babel}
\usepackage{amsmath}
\usepackage{amssymb}
\usepackage{array}
\usepackage{booktabs}
\usepackage{caption}
\usepackage{subcaption}
\usepackage{enumitem}
\usepackage{fancyhdr}
\usepackage{titlesec}
\usepackage{url}
\usepackage{textcomp}
\usepackage{xcolor}
\usepackage{listings}
\usepackage{multirow}
\usepackage{longtable}
\usepackage{tabularx}
\usepackage{times}

% Set up page style
\pagestyle{fancy}
\fancyhf{}
\rfoot{\thepage}
\renewcommand{\headrulewidth}{0pt}

% Title formatting
\titleformat{\chapter}[display]{\normalfont\bfseries}{\chaptertitle\thechapter}{0pt}{\Huge}
\titlespacing*{\chapter}{0pt}{50pt}{40pt}

% Document info
\title{Template Document}
\author{Academic Support Agent}
\date{\today}

\begin{document}
\maketitle

% Content from titlepage would be included here

\pagenumbering{roman}

\newpage

\newpage

\newpage
\titleformat{\section}
  {\normalfont\Large\bfseries}{\thesection}{1em}{}[{\color{brightmaroon}\titlerule[0.8pt]}]
\pagenumbering{arabic}

\section{Executive summary}

Do not forget that this needs to be included. Check the thesis handbook for details!!!


\section{This could be your introduction}
Text with citation according to~\cite{Ref1}.

\subsection{Is your thesis interesting and relevant?}
Text

\subsection{What questions do you address and how?}
Text

\subsection{How is it structured?}
Text


\section{A next section}
\subsection{With one subsection}
All of this is summarized in Figure~\ref{fig:LSCM}.

\begin{figure}[htbp!]
\centering
\includegraphics[width=0.3\linewidth]{images/LSCM.jpg}
%\caption{case 3: combined situation}
\caption{Features of the LSCM study program}
\label{fig:LSCM}
\end{figure}

\subsection{And another subsection}
Table~\ref{tab:table1} is a cool table. Avoid too fancy formatting and focus on clear, easy to read and focused information in a table. Do not forget that you may use appendices to add extensive information that exceeds the limit one wants to put in a table in the body of text. 

\begin{table}[htp!]
\centering
\caption{A clean table is a nice thing}
\label{tab:table1}
\begin{tabular}{@{}ccc@{}}
\toprule
First Column & Second Column & Third Column \\ \midrule
Text & Text & Text \\
Text & Text & Text \\
Text & Text & Text \\ \bottomrule
\end{tabular}
\end{table}

\section{Yet another section}


\section{Conclude your thesis with conclusions}


\newpage
%% The Appendices part is started with the command 
%% appendix sections are then done as normal sections
%\section*{Acknowledgements}



%% References
%%
%% Following citation commands can be used in the body text:
%% Usage of \cite is as follows:
%%   \cite{key}         ==>>  [#]
%%   \cite[chap. 2]{key} ==>> [#, chap. 2]
%%

%% References with bibTeX database:

%

% 
% 
% 
% 
%
% 
% 
% 
%
% 


It is strongly advised that you use special software support to manage your list of references, such as e.g. EndNote, Citavi or Mendeley. In this Overleaf template, don't forget to use the sample.bib  file.

% Appendices for the result tables. Appendix A is dedicated to the base case results, and Appendix B to the parametric analysis results.
\newpage

\begin{appendices}

% Appendix A
\section{Apendix Title}

% Appendix B
\section{Apendix Title}

\end{appendices}
\end{document}